\documentclass[tikz,border=6pt]{standalone}
\usepackage{peabra}
\begin{document}
\begin{tikzpicture}[x=1cm,y=1cm, a/.style={doboz=29mm, font=\footnotesize, minimum height=14mm}]
  \foreach \n/\t/\x/\f in {m/{küldetés- és forgatókönyv-elemzés}/0/{4. fejezet},
                           fe/{funkcióelemzés}/3.3/{8. fejezet},
                           fm/{funkciómegosztás ember és gép között}/6.6/{12. fejezet},
                           fa/{feladatelemzés}/9.9/{11. fejezet},
                           rt/{rendszertervezés}/13.2/{8--10. fejezet},
                           te/{teszt és értékelés}/16.5/{23--24. fejezet}} {
    \node[a] (\n) at (\x,0) {\t};
    \node[felirat] at (\x,-1.05) {\f};
  }
  \foreach \a/\b in {m/fe,fe/fm,fm/fa,fa/rt,rt/te} \draw[nyil] (\a) -- (\b);
  \draw[nyil, draw=peKiemel, dash pattern=on 3pt off 2pt] (te.north) -- ++(0,0.75) -| node[felirat, pos=0.25, above, text=peKiemel] {a felhasználói értékelés eredményei visszahatnak a megosztásra és a feladatokra} (fm.north);
  \node[font=\cimbetu\small, text=peKek, anchor=west] at (-1.45,2.45) {ELEMZÉS};
  \node[font=\cimbetu\small, text=peKiemel, anchor=east] at (17.95,2.45) {TERVEZÉS, TESZT, ÉRTÉKELÉS};
\end{tikzpicture}
\end{document}
